%% notes-tex/030-scaling-plots/sections/1-inputs.tex

\section{What the runs provide\secstatus{todo}}
\label{sec:inputs}

\begin{keybox}
\textbf{Measured.} The rows each run logs and the reads the readout takes from
them are implemented and tested (\texttt{scaling\_readout\_030.py},
\texttt{torchcell.trainers.compute\_accounting}). The grid is queued on Delta.

\textbf{Not measured.} Every curve in every sketch. No run has reported, and no
plot here has a data point behind it.
\end{keybox}

The grid, its reads and its limits are set out in the companion study
(\texttt{notes-tex/030-scaling-study}), whose vocabulary applies here. In short:
the \term{loss} is the validation mean squared error on the gene-interaction
score in \term{label units}, after the inverse of the training normalization,
pooled over the pinned 37,705 triples and per \term{screen} (Kuzmin 2018,
Kuzmin 2020) through the per-token keys; $N$ is non-embedding parameters; $D$
is distinct training records, and on 030 also gene-interaction entry rows; the
\term{selected read} is the loss at the epoch of lowest pooled validation MSE;
the \term{matched read} is the loss interpolated at one shared records-seen
budget $B$, the smallest total any point of the grid reached; the
\term{replicate floor} $E = (1 - r)\,\mathrm{var}(\tau)$ is the loss no predictor
of a single measurement can beat, per screen; the \term{held-out test} refits
without the largest point of an axis and predicts it.

\notesec{Rows}
Each run logs, at the end of every epoch, the pooled and per-screen validation
MSE and Pearson, the records seen, the cumulative measured FLOPs and GPU hours,
and the epoch index that joins them. A point is one (size, subset, seed, learning
rate); a chain of jobs is one point with its epochs stitched by the checkpoint
each link resumed from. The readout keeps every epoch of every point
(\texttt{<sweep>\_epochs.csv}), one row per point and read
(\texttt{<sweep>\_points.csv}), every fit with its interval and held-out
prediction (\texttt{<sweep>\_fits.json}), and the gap per subset
(\texttt{<sweep>\_gap.csv}). Every plot below is drawn from one of those four
files, so a plot is regenerated from the rows and never assembled by hand.

\notesec{The catalog}
Table~\ref{tab:catalog} is the list. The id is the panel letter family used in
the sketches: S for the loss surface, C for compute, D for diagnostics, P for the
pair. Status is read against the queued grid as it stands: full means every row
the plot needs exists at the design's seed count; partial means the rows exist
but a limit of the design narrows what the plot can claim, named in its section.

%% GENERATED FILE -- do not hand-edit.
%% SOURCE: experiments/030-solid-growth-multi/scripts/scaling_plot_catalog_030.py (the PlotSpec catalog)
\begin{table}[htbp]
\centering\footnotesize
\caption{The plot catalog: 16 plots the queued grid can feed, 9 in full and 7 in part. Read names what the y value is at each point: per-epoch rows, the validation-selected epoch, the interpolation at the shared records-seen budget $B$, or the calibration jobs. Feeds names the rows of the grid the plot needs. Partial means the rows exist but a stated limit narrows the claim (ten-epoch full-pool points, one seed, four points on an axis).}
\label{tab:catalog}
\begin{tabular}{llp{0.30\linewidth}p{0.16\linewidth}p{0.26\linewidth}l}
\toprule
id & family & plot (x; y) & read & feeds & status \\
\midrule
S1 & surface & \textbf{Loss against D at every size}: D, training records (log); 030 also in entry rows; val MSE, label units (log) & both & the 5 x 4 seed-1 grid; seeds 2-3 at the two smallest subsets for 45x2 and 180x8 & full \\
S2 & surface & \textbf{Loss against N at every subset}: N, non-embedding parameters (log); val MSE, label units (log) & both & the 5 x 4 seed-1 grid & partial \\
S3 & surface & \textbf{The joint surface as a 5 x 4 heatmap}: D, four subsets; N, five sizes & both & the 5 x 4 seed-1 grid & full \\
S4 & surface & \textbf{Fitted forms on the two axes}: D or N (log); val MSE minus the fitted E (log) & both & S1's 180x8 column and S2's full-pool column & partial \\
C1 & compute & \textbf{Learning curves against records seen}: records seen (log); val MSE, label units (log) & per-epoch & every run; chains stitched by resumed\_from & full \\
C2 & compute & \textbf{Loss against training FLOPs, the envelope}: cumulative FLOPs (log); val MSE, label units (log) & per-epoch & every run on the full pool and on 1/4 (the large-data columns) & partial \\
C3 & compute & \textbf{Loss against GPU hours}: GPU hours (log); val MSE, label units (log) & per-epoch & every run & full \\
C4 & compute & \textbf{Cost per record and pace against size}: N (log); FLOPs per record (left); records per second (right) & calibration & the two calibration jobs (180x8, 45x2) now; every grid run's first epoch after & partial \\
D1 & diagnostic & \textbf{Bootstrap distribution of each exponent}: exponent value; count of resamples & both & 300 resamples over runs of S1's and S2's fitted columns & full \\
D2 & diagnostic & \textbf{Held-out prediction against observation}: predicted loss at the held-out point; observed loss & both & every fit of S4 & partial \\
D3 & diagnostic & \textbf{Fitted E against the replicate floor}: screen (2018, 2020, pooled); loss, label units & both & S4's floored fits and the noise-ceiling file & full \\
D4 & diagnostic & \textbf{Seed spread against effect size}: (size, subset) cell; val MSE, label units & selected & the eight seed-replicate cells plus their seed-1 neighbors & full \\
D5 & diagnostic & \textbf{The learning-rate check}: epoch; val MSE, label units (log) & per-epoch & the two lr jobs and the seed-1 cell at 45x2 on 1/4 & partial \\
D6 & diagnostic & \textbf{Per-screen Pearson beside the loss}: D (log); val Pearson, per screen and pooled & selected & the 180x8 column & full \\
P1 & pair & \textbf{The 030 minus 025 gap against D}: D, training records (log); 030 minus 025 pooled val MSE, label units & both & the 180x8 column of both sweeps & partial \\
P2 & pair & \textbf{Paired learning curves}: records seen (log); val MSE, label units (log) & per-epoch & the 180x8 column of both sweeps & full \\
\bottomrule
\end{tabular}
\end{table}

